\section{Related Work}

Our work builds on three research areas: SMT-guided program repair for human code, neural code generation without formal verification, and static analysis for typed languages. We position our hypothesis---incremental SMT for LLM code repair in typed Python---at the intersection, though empirical validation was blocked by infrastructure failures.

\subsection{SMT-Guided Program Repair}

\textbf{Angelix and Prophet} apply SMT solvers to synthesize patches for buggy programs by encoding correctness specifications as constraints. These systems work on small human-written patches (typically 1-10 lines) and assume localized modifications. Their constraint extraction relies on predefined templates for common bug patterns (off-by-one errors, missing null checks). While effective for human code repair, they do not scale to full LLM-generated programs where bug locations are unknown and programs exceed 100 lines of code.

\textbf{Our hypothesis} extends this paradigm to LLM code repair workflows in typed Python: instead of batch re-verification on each iteration, incremental SMT re-verifies only modified functions and dependencies. The key difference is scale---we target full programs, not patches---and code source---LLM-generated code may have different quality characteristics than human code. However, this transferability assumption (A1: LLM code has extractable constraints) remains untested due to our validation failure.

\subsection{Neural Code Generation}

\textbf{AlphaCode} and \textbf{CodeT5} demonstrate that large-scale pretraining enables LLMs to generate functionally correct code on benchmarks like Codeforces and HumanEval. These systems rely on test-suite validation: generate multiple candidates, filter by passing tests, rank by confidence. While effective for competition programming, test suites provide incomplete correctness guarantees---they miss edge cases not covered by tests.

\textbf{Formal verification} could address this gap by providing SMT-based correctness proofs. However, existing neural code generation systems do not integrate static analysis or SMT solvers. We proposed adding incremental SMT verification to LLM repair loops, but we could not validate whether LLM-generated code has extractable constraints (assumption A1 untested due to infrastructure failure) or whether SMT provides practical benefit over test-suites (comparison untested).

\subsection{Static Analysis for Typed Languages}

\textbf{Prusti} (Rust) and \textbf{Pyre} (Python) extract SMT constraints from type annotations and function contracts (preconditions, postconditions, invariants). These tools demonstrate that type annotations map to first-order logic predicates suitable for SMT solving. Prusti uses Rust's ownership types and developer-written specifications; Pyre uses Python type hints and optional runtime assertions.

\textbf{Transferability to LLM Code:} Prior work validates these analyzers on human-written typed code. Whether LLM-generated code has sufficient annotation quality for constraint extraction is an open question. We hypothesized $\geq$90\% extraction success (assumption A1), but we could not test this due to API authentication failures blocking code generation. The constraint extraction pipeline was validated on error files (finding 0 constraints as expected), confirming architectural soundness on negative cases, but the core claim---that LLM code is analyzable---remains empirically unverified.

\subsection{Neural Program Repair}

\textbf{CURE} and \textbf{CoCoNut} demonstrate that neural program repair typically modifies 1-3 lines in human code (80\%+ of repairs are localized). Our hypothesis assumes LLM repairs exhibit similar locality (assumption A2, untested). If true, dependency cones remain small (<30\%), enabling incremental SMT speedup. However, LLM repair patterns may differ from human edits---broader modifications, cross-module dependencies---which could invalidate speedup claims. We could not test repair locality due to prerequisite failures (h-m2, h-m3 blocked). Assumption A2 transfers from human code literature without validation; LLM repair patterns may differ significantly, potentially eliminating speedup benefits even if extraction succeeds.

\subsection{Benchmarks for Typed Code}

\textbf{HumanEval} provides 164 Python programming problems with test suites but lacks type annotations required for static analysis. \textbf{APPS} and \textbf{CodeContests} similarly focus on functional correctness without formal specifications. Our contribution---HumanEval + Pydantic extensions---fills this gap: we mechanically generated 100 type-annotated prompts. This artifact is reusable for future constraint extraction research, independent of our validation failure.

\subsection{Positioning}

Our hypothesis occupies unexplored territory: combining SMT-guided repair (established for human patches) with neural code generation (lacking verification) in typed Python with Pydantic. The key novelty is incremental SMT for full LLM programs, but this remains a proposal pending validation. The verified contribution is the dataset artifact and pipeline architecture, rather than performance claims (speedup untested).
